% ============================================================
% Chapter 1: Introduction
% ============================================================
\chapter{Introduction}
\label{ch:introduction}

\chapteroverview{This chapter introduces the fundamental clinical background, problem statement, and engineering objectives of the proposed EOG-based assistive technology system. It contextualizes non-invasive Human-Machine Interfaces (HMIs) for patients with Locked-in Syndrome (LIS), outlines the specific project goals, presents the structure of the thesis report, and details the two-year engineering milestone execution timeline.}


% ============================================================
% Section 1.1: General Overview and Clinical Background
% ============================================================
\section{General Overview and Clinical Background}
\label{sec:general_overview_clinical_background}

Severe neurodegenerative disorders and traumatic injuries---such as advanced Amyotrophic Lateral Sclerosis (ALS), Muscular Dystrophy, brainstem stroke, and high-level spinal cord injuries---frequently result in profound motor paralysis, often presenting as Locked-in Syndrome (LIS) or Quadriplegia. Patients in these clinical conditions suffer from a complete loss of voluntary muscle control throughout the body, depriving them of conventional speech and physical interaction with their environment. However, the cranial nerves controlling extraocular muscle motility remain physiologically intact and functionally operational in most affected individuals.

To address these severe functional deficits, non-invasive Human-Machine Interfaces (HMIs) based on electrooculography (EOG) have emerged as an effective assistive solution. EOG exploits the standing electrostatic potential maintained between the cornea and the retina to track eye displacement and voluntary eyelid blinks.


% ============================================================
% Section 1.2: Problem Statement
% ============================================================
\section{Problem Statement}
\label{sec:problem_statement}

Individuals with severe motor paralysis face extreme social isolation and total reliance on caregivers due to the lack of functional communication and mobility tools. Although various assistive technologies currently exist, they present major technical and practical limitations. First, invasive Brain-Computer Interfaces (BCIs) provide direct neural control but carry significant surgical risks, tissue infection hazards, and high financial costs. Second, Video-Oculography (VOG) camera-based eye-tracking alternatives are highly sensitive to ambient lighting variations, require continuous spatial calibration, and demand high computational processing power. Third, system fragmentation remains a major barrier, as most commercially available devices treat communication and physical locomotion as separate challenges, requiring independent and expensive standalone platforms for each function.


Consequently, an engineering need exists for a non-invasive, low-cost HMI with reliable threshold detection. Such a system must utilize simplified biopotential acquisition hardware to deliver real-time control over both digital text communication and motorized mobility platforms within a single interface.

This project presents the design and implementation of an assistive technology system that integrates a virtual keyboard for text-based communication and a prototype motorized wheelchair controlled via real-time EOG thresholding. This framework addresses mobility and communication requirements within a single, unified platform.


% ============================================================
% Section 1.3: Project Engineering Objectives
% ============================================================
\section{Project Engineering Objectives}
\label{sec:project_engineering_objectives}

The primary objective of this project is to develop an integrated, cost-effective, non-invasive HMI that provides individuals with locked-in syndrome (LIS) an assistive interface for communication and navigation. To attain this overall goal, the following specific engineering objectives were established:

\begin{enumerate}[label=\textbf{\arabic*.}]
    \item \textbf{Bio-Potential Acquisition Hardware Design:} To design and assemble a high-gain, low-noise analog front-end circuit capable of reliably capturing microvolt-level corneo-retinal potentials while rejecting power-line interference and environmental electromagnetic noise.
    \item \textbf{Real-Time Signal Processing Framework:} To implement real-time digital filtering and thresholding algorithms within a hybrid MATLAB-Arduino processing architecture to extract and classify the temporal characteristics of voluntary eye blinks.
    \item \textbf{Finite State Machine Control Integration:} To engineer a deterministic Finite State Machine (FSM) that translates multi-blink sequences into operational commands, allowing a single biopotential channel to manage two distinct subsystems while enforcing safety protocols.
    \item \textbf{Virtual Optical Speller Development:} To construct an auto-scanning virtual keyboard featuring a Graphical User Interface (GUI) configured to reduce cognitive load, enabling users to compose, edit, and clear text.
    \item \textbf{3D-Printed Prototype Wheelchair Implementation:} To construct a functional 3D-printed motorized wheelchair prototype that receives wireless directional commands, serving as a scaled experimental model to validate wheelchair navigation and obstacle-avoidance logic.
\end{enumerate}

% ============================================================
% Section 1.4: Report Organization
% ============================================================
\section{Report Organization}
\label{sec:report_organization}

This thesis is organized into six main chapters structured logically to document the complete engineering lifecycle of the EOG-based assistive system:

\begin{itemize}
    \item \textbf{Chapter 1 (Introduction):} Outlines the clinical background, problem statement, engineering objectives, thesis structure, and execution timeline.
    \item \textbf{Chapter 2 (Literature Review and Theoretical Background):} Reviews existing literature on biopotential acquisition, HMIs, virtual spellers, and eye-tracking systems while identifying current technical gaps.
    \item \textbf{Chapter 3 (System Architecture and Block Diagram):} Details the high-level architecture, functional block diagrams, operational mode-switching logic, and overall system flowcharts.
    \item \textbf{Chapter 4 (Circuit Simulation and Hardware Implementation):} Presents circuit simulations, mathematical design calculations, component selection (AD620, L298N), and software architecture.
    \item \textbf{Chapter 5 (Results and Discussion):} Evaluates experimental biopotential waveforms, obstacle avoidance response times, virtual speller accuracy, and statistical performance metrics.
    \item \textbf{Chapter 6 (Conclusion and Recommendations):} Synthesizes the core engineering findings, clinical applications, and provides strategic recommendations for future technical developments.
\end{itemize}

% ============================================================
% Section 1.5: Project Execution Timeline
% ============================================================
\section{Project Execution Timeline}
\label{sec:project_execution_timeline}

The development lifecycle of this project was systematically organized across two academic years (within the four-year Biomedical Engineering curriculum).

During Year 3, efforts focused on EOG signal acquisition, continuous hardware design, noise suppression, and circuit optimization until a stable signal output was achieved. Building upon this hardware foundation, Year 4 focused on software development, virtual keyboard GUI programming, finite state machine integration, 3D printing of the prototype vehicle, and wireless control integration. Table \ref{tab:project_timeline} summarizes the engineering milestones achieved across this timeline.

\begin{table}[H]
  \centering
  \caption{Two-Year Project Execution Timeline (Year 3 -- Year 4)}
  \label{tab:project_timeline}
  \begin{tabularx}{\textwidth}{l X c}
    \toprule
    \textbf{Stage} & \textbf{Engineering Task \& Technical Focus} & \textbf{Academic Period} \\
    \midrule
    Phase 1 & Clinical background research, initial EOG signal acquisition, and preliminary analog circuit prototyping & Year 3 -- Semester 1 \\
    Phase 2 & Hardware optimization, active filter tuning, and achieving stable signal acquisition & Year 3 -- Semester 2 \\
    Phase 3 & Virtual keyboard programming (GUI), FSM control logic implementation, and signal algorithm refinement & Year 4 -- Semester 1 \\
    Phase 4 & 3D-printed vehicle construction, wireless control integration, system validation testing, and final documentation & Year 4 -- Semester 2 \\
    \bottomrule
  \end{tabularx}
\end{table}


% --- Chapter Conclusion & Transition ---
\vspace{1em}
In summary, this introductory chapter established the clinical motivation, engineering scope, and milestone timeline of the proposed EOG-guided assistive system. Having outlined the foundational objectives for vehicle mobility and virtual speller communication, Chapter 2 reviews relevant literature, the electrophysiology of the corneo-retinal dipole, and biopotential conditioning principles.
